\section*{الفصل \ref{ch:b2:maxwell-equations} --- معادلات ماكسويل}

\begin{solution}{exo:b2:maxwell-equations:1}
(أ) $\operatorname{div} = 3$، و $\operatorname{\vect{curl}} = \vect 0$: أي \omterm{def:b2:maxwell-equations:operators}{تدرج} $r^2/2$.
(ب) $0$ و $(0, 0, 2)$: أي \omterm{def:b2:maxwell-equations:operators}{دوّار} (للمجال $-\tfrac12(x^2 + y^2)\vect e_z$)، لا
تدرجًا. (ج) $0$ و $\vect 0$: أي \omterm{def:b2:maxwell-equations:operators}{تدرج} $xyz$ (\omterm{def:b2:maxwell-equations:operators}{ودوّار} كذلك). (د)
$\frac1{r^2}\partial_r(r^2/r^2) = 0$، \omterm{def:b2:maxwell-equations:operators}{والدوّار} $\vect 0$: أي \omterm{def:b2:maxwell-equations:operators}{تدرج} $-1/r$. (ه)
$0$ و $(0, -2x, 0)$: أي \omterm{def:b2:maxwell-equations:operators}{دوّار}، لا تدرجًا.
\end{solution}

\begin{solution}{exo:b2:maxwell-equations:2}
(أ) غاوس وفاراداي: طرفاهما بواحدة \unit{V/m^2} (أي كثافة شحنة على
$\varepsilon_0$، أو \omterm{def:b2:maxwell-equations:operators}{تدرج} مجال، أو مجال على زمن)؛ وأمبير: طرفاه بواحدة
\unit{T/m}.
(ب) $1/\sqrt{8.85 \times 10^{-12} \times 1.257 \times 10^{-6}} = \qty{3.00e8}{m/s}$. (ج) والنسبة
$\varepsilon_0\omega/\gamma$: النحاس $5 \times 10^{-20}$؛ وماء البحر عند \qty{1}{GHz} $1 \times 10^{-2}$
(وهو ناقل بالكاد)؛ والزجاج عند \qty{50}{Hz} $3 \times 10^3$: أي عازل،
ويهيمن فيه \omterm{thm:b2:maxwell-equations:maxwell}{تيار الإزاحة}.
\end{solution}

\begin{solution}{exo:b2:maxwell-equations:3}
$\dd E/\dd t = I/\varepsilon_0\pi R^2 = \qty{1.8e12}{V.m^{-1}.s^{-1}}$؛ و $j_d = I/\pi R^2 =
\qty{16}{A/m^2}$؛ و $B(\qty{5}{cm}) = \mu_0Ir/2\pi R^2 = \qty{0.5}{\micro T}$؛ و $B(\qty{20}{cm}) =
\mu_0I/2\pi r = \qty{0.5}{\micro T}$ --- أي مجال السلك بالضبط.
\end{solution}

\begin{solution}{exo:b2:maxwell-equations:4}
(أ) $\rho = \varepsilon_0k$. (ب) $\frac1{r^2}\partial_r(\rho_0r^3/3\varepsilon_0) = \rho_0/\varepsilon_0$؛
وفي الخارج، بغاوس: $E = \rho_0R^3/3\varepsilon_0r^2$. (ج) $\operatorname{div}\vect B = 0$: نعم؛
و $\operatorname{\vect{curl}}\vect B = -(B_0/a)\vect e_z$: ومنه $\vect j = -(B_0/\mu_0a)\vect e_z$، أي
صفيحة تيار منتظمة سماكتها على امتداد $y$.
\end{solution}

\begin{solution}{exo:b2:maxwell-equations:5}
(أ) $\vect A = \tfrac12(B_yz - B_zy, B_zx - B_xz, B_xy - B_yx)$: ومنه $(\operatorname{\vect{curl}}
\vect A)_x = \tfrac12(B_x + B_x) = B_x$، وهكذا؛ و $\operatorname{div}\vect A = 0$. (ب)
$\vect A' - \vect A = \tfrac12B_0(y, x, 0) = \operatorname{\vect{grad}}(\tfrac12B_0xy)$. (ج)
$\frac1r\partial_r(Br^2/2) = B$؛ و $\frac1r\partial_r(BR^2/2) = 0$؛ وكلاهما يساوي $BR/2$ عند
$R$. (د) $2\pi r\cdot BR^2/2r = \pi R^2B = \Phi$: فحلقة خارج الملف اللولبي،
حيث $\vect B = \vect 0$، تحسّ مع ذلك $e = -\dd\Phi/\dd t$ --- عبر $\vect E =
-\partial_t\vect A$، وهو لا ينعدم هناك.
\end{solution}

\begin{solution}{exo:b2:maxwell-equations:6}
(أ) $E = \sigma/\varepsilon_0 = \qty{1.1e5}{V/m}$، ناظمي. (ب) $\sigma/\varepsilon_0$ بينهما، ومعدوم
في الخارج (فكل قفزة تضيف $\pm\sigma/\varepsilon_0$). (ج) $j_s = nI = \qty{1e4}{A/m}$، والقفزة
$\mu_0j_s = \qty{12.6}{mT} = B$ في الداخل. (د) وحلقة خارج الحلقة الطارية لا تحيط
بأي تيار صافٍ، ومنه $B = 0$ هناك؛ وفي الداخل يعطي أمبير $\mu_0NI/2\pi r$،
وتتوافق معه القفزة عبر اللفّ، $\mu_0j_s = \mu_0NI/2\pi r$.
\end{solution}

\begin{solution}{exo:b2:maxwell-equations:7}
(أ) $V'' = -\rho_0/\varepsilon_0$: ومنه $V = -\rho_0x^2/2\varepsilon_0 + (U/d + \rho_0d/2\varepsilon_0)x$. (ب)
$E = \rho_0x/\varepsilon_0 - U/d - \rho_0d/2\varepsilon_0$، وينعدم عند $x_0 = d/2 + \varepsilon_0U/\rho_0d$
--- أي بجوار المنتصف من أجل شحنة فراغية كبيرة. (ج) $\sigma(0) = \varepsilon_0E(0^+) =
-\varepsilon_0U/d - \rho_0d/2$، و $\sigma(d) = -\varepsilon_0E(d^-) = \varepsilon_0U/d - \rho_0d/2$: ومجموعهما
$-\rho_0d$، فيحيّد الشحنة الفراغية. (د) ومع $\rho_0 = 0$: $V$ خطي و
$E$ منتظم؛ ومع $U = 0$: قطع مكافئ ذروته عند $d/2$.
\end{solution}

\begin{solution}{exo:b2:maxwell-equations:8}
(أ) \qty{6000}{km}، و \qty{300}{m}، و \qty{3}{m}، و \qty{12.5}{cm}. (ب) $f < c/10L$:
أي \qty{300}{MHz}؛ و \qty{3}{GHz}؛ و \qty{30}{Hz}. (ج) $(2\pi \times 10^6 \times 0.05/3 \times 10^8)^2/4 =
3 \times 10^{-7}$. (د) وهي ليست كذلك بدقة: فعند \qty{50}{Hz} تكون شبكة طولها \qty{1000}{km}
سدس طول موجة، وتختلف الأطوار عبرها، وتُعامَل الخطوط الطويلة
خطوطَ نقل --- غير أن كل دارة محلية
شبه مستقرة.
\end{solution}

\begin{solution}{exo:b2:maxwell-equations:9}
كان $\operatorname{div}\operatorname{\vect{curl}}\vect B = 0$ ليفرض $\operatorname{div}\vect j =
0$، أي $\partial_t\rho = 0$ في كل مكان. وفي اللبوس (وحجمه $\pi R^2e$) يدخل
التيار $I$ ولا يخرج شيء: ومنه $\langle\operatorname{div}\vect j\rangle =
-I/\pi R^2e = \qty{-1.6e5}{A/m^3}$، و $\partial_t\rho = \qty{+1.6e5}{C.m^{-3}.s^{-1}}$.
\end{solution}

\begin{solution}{exo:b2:maxwell-equations:10}
(أ) $2\pi rE_\theta = -\pi r^2\dot B$: ومنه $E_\theta = -r\dot B/2$ في الداخل، و $-R^2\dot B/2r$
في الخارج. (ب) وخطوطه دوائر مغلقة: فهو ليس تدرجًا، بل
يقود الشحن حول حلقة --- أي مجال محرّك كهربائيًّا. (ج) $e\cdot2\pi r_0
|E_\theta| = e\pi r_0^2\dot B$: أي $\pi \times 0.25 \times 10 = \qty{7.9}{eV}$ في الدورة؛ و $2.5 \times 10^6$
دورة. (د) ويعطي $p = eB(r_0)r_0$ العلاقة $\dot p = er_0\dot B(r_0)$، بينما يعطي المجال
المحرَّض $\dot p = e|E_\theta| = e\dot\Phi/2\pi r_0$: ومنه $B(r_0) = \Phi/2\pi r_0^2 =
\tfrac12\langle B\rangle$.
\end{solution}

\begin{solution}{exo:b2:maxwell-equations:11}
(أ) $\operatorname{\vect{curl}}\operatorname{\vect{curl}}\vect E = -\Delta\vect E$ (لأن $\operatorname{div}
\vect E = 0$) $= -\partial_t\operatorname{\vect{curl}}\vect B = -\mu_0\varepsilon_0\partial_t^2\vect E$. (ب)
والأمر نفسه بتبادل الدورين. (ج) $c = 1/\sqrt{\mu_0\varepsilon_0} = \qty{3.0e8}{m/s}$.
(د) كان فيبر وكولراوش قد قاسا نسبة واحدتي الشحنة
الكهرسكونية والكهرمغناطيسية، وهي $3.1 \times 10^8$ متر في الثانية؛ وكان فيزو قد
قاس الضوء عند $3.15 \times 10^8$ متر في الثانية: «فلا نكاد نستطيع تجنّب
الاستنتاج بأن الضوء تموّجات مستعرضة في
الوسط نفسه».
\end{solution}

\begin{solution}{exo:b2:maxwell-equations:12}
(أ) $\operatorname{\vect{curl}}\operatorname{\vect{curl}}\vect B = -\Delta\vect B = \mu_0\gamma\operatorname{\vect{curl}}
\vect E = -\mu_0\gamma\partial_t\vect B$: ومنه $\Delta\vect B = \mu_0\gamma\partial_t\vect B$، وهي صيغة
معادلة الحرارة بانتشارية $1/\mu_0\gamma$. (ب) $\tau \sim \mu_0\gamma d^2$: أي \qty{7.5}{ms}
للصفيحة؛ و $4\pi \times 10^{-7} \times 5 \times 10^5 \times 9 \times 10^{12} = \qty{6e12}{s}$، أي مئتا
ألف سنة، للّبّ. (ج) $\mu_0\gamma\delta^2 \sim 2/\omega$: ومنه $\delta = \sqrt{2/\mu_0
\gamma\omega}$: أي \qty{9}{mm} عند \qty{50}{Hz}، و \qty{65}{\micro m} عند \qty{1}{MHz}. (د) ولا يستطيع مجال
اللبّ أن يتغير بالانتشار في أقل من $10^5$ سنة: فالتغيّر
الدهري المرصود والانقلابات تقتضي حركة المائع
الناقل --- أي مولّدًا ذاتيًّا.
\end{solution}

\begin{solution}{pb:b2:maxwell-equations:1}
\textbf{1.} $Q = (I_0/\omega)\sin\omega t$، و $E = Q/\varepsilon_0\pi R^2$: وسعته $I_0/\omega\varepsilon_0
\pi R^2 = \qty{5.7e5}{V/m}$.

\textbf{2.} $j_d = \varepsilon_0\partial_tE = I(t)/\pi R^2 = 32\cos\omega t$ A/m$^2$؛ وتدفقه
عبر اللبوس هو $I(t)$.

\textbf{3.} $2\pi rB = \mu_0\varepsilon_0\pi r^2\partial_tE$: ومنه $B = \mu_0I(t)r/2\pi R^2$؛ وعند $R$:
$\mu_0I_0/2\pi R = \qty{2}{\micro T}$.

\textbf{4.} من أجل $r > R$ يكون \omterm{thm:b2:maxwell-equations:maxwell}{تيار الإزاحة} كله محاطًا: ومنه $B =
\mu_0I/2\pi r$، أي مجال السلك؛ وهو متصل عند $R$.

\textbf{5.} $(\operatorname{\vect{curl}}\vect E_1)_\theta = -\partial_rE_1 = -\partial_tB_\theta$، ومنه $\partial_rE_1
= \mu_0\varepsilon_0(r/2)\ddot E$ و $E_1(r) - E_1(0) = \mu_0\varepsilon_0r^2\ddot E/4$؛ ومع
$\ddot E = -\omega^2E$: $|E_1(R) - E_1(0)|/|E| = (\omega R/c)^2/4 = 1 \times 10^{-6}$.

\textbf{6.} ممتازة. و $10\%$ تقتضي $(\omega R/c)^2 = 0.4$، أي $f \approx \qty{300}{MHz}$ ---
فيصير جوفًا، لا مكثفًا.

\textbf{7.} $W_e = \tfrac12\varepsilon_0E_0^2\pi R^2d = \qty{4.5e-5}{J}$؛ و $W_m = \int_0^R(\mu_0I_0r/
2\pi R^2)^2/2\mu_0\cdot2\pi rd\,\dd r = \mu_0I_0^2d/16\pi = \qty{2.5e-11}{J}$؛ والنسبة $6 \times
10^{-7} \approx (\omega R/c)^2/8$: أي الوسيط الصغير نفسه.

\textbf{8.} $\oint\vect E\cdot\dd\vect l = 2\pi r_0E_\theta = -\dd\Phi/\dd t$.

\textbf{9.} $\dd p/\dd t = e|E_\theta| = (e/2\pi r_0)|\dd\Phi/\dd t|$؛ وفي الدورة،
$2\pi r_0\cdot e|E_\theta| = e|\dd\Phi/\dd t|$.

\textbf{10.} يعطي $p = eB(r_0)r_0$ العلاقة $\dot p = er_0\dot B(r_0)$؛ وبمساواتها بالسؤال 9:
$B(r_0) = \Phi/2\pi r_0^2 = \tfrac12\langle B\rangle$.

\textbf{11.} في المجال المنتظم يكون $B(r_0) = \langle B\rangle$، أي ضعف ما يلزم:
فتنمو القوة المغناطيسية أسرع من كمية الحركة ويتقلّص
المدار. ويُشكَّل القطبان بحيث يكون المجال قويًّا بجوار
المحور ونصف ذلك على المدار.

\textbf{12.} $p = eBr_0 = \qty{3.2e-20}{kg.m/s}$؛ و $E = pc = \qty{9.6e-12}{J} = \qty{60}{MeV}$؛
و $\Phi = 2\pi r_0^2B = \qty{0.63}{Wb}$ في \qty{5}{ms}: أي \qty{126}{eV} في الدورة، و $4.8 \times
10^5$ دورة، و \qty{1500}{km} --- وبسرعة $c$ في \qty{5}{ms}: وهو متوافق.

\textbf{13.} من مصدر تغذية المغنطيس، مسلَّمةً إلى الإلكترونات عبر
حدّ فاراداي: فالتدفق المتغير يصنع المجال المسرِّع.

\textbf{14.} عشرات الميغاإلكترون فولت من الإلكترونات تكبح في هدف ثقيل فتعطي
أشعة سينية نافذة؛ وقد حلّ محلّه المسرّع الخطي المدمج، وهو أعلى معدلًا
في الجرعة.

\textbf{15.} $\operatorname{div}\operatorname{\vect{curl}}\vect A = 0$؛ و $\operatorname{\vect{curl}}(-\operatorname{\vect{grad}}
V - \partial_t\vect A) = -\partial_t\operatorname{\vect{curl}}\vect A = -\partial_t\vect B$.

\textbf{16.} $\frac1{r^2}(r^2V')' = -\rho_0/\varepsilon_0$: ومنه $V = \rho_0(3a^2 - r^2)/6\varepsilon_0$
في الداخل، و $\rho_0a^3/3\varepsilon_0r = Q/4\pi\varepsilon_0r$ في الخارج؛ و $E = \rho_0r/3\varepsilon_0$ و
$Q/4\pi\varepsilon_0r^2$.

\textbf{17.} $\tfrac12\int\rho V\,\dd\tau = 2\pi\rho_0^2(a^5/6 - a^5/30)/\varepsilon_0 = 4\pi\rho_0^2a^5/
15\varepsilon_0 = 3Q^2/20\pi\varepsilon_0a$؛ واليورانيوم: \qty{1.6e-10}{J} $\approx \qty{1000}{MeV}$ (والانشطار
يحرّر \qty{200}{MeV} التي تكون بها كرتان أصغر أرخص).

\textbf{18.} $V = -E(t)z$ (مضافًا إليه ثابت). ومع $\vect A = A_z(r, t)\vect e_z$،
يعطي $B_\theta = -\partial_rA_z = \mu_0Ir/2\pi R^2$ القيمة $A_z = -\mu_0I(t)r^2/4\pi R^2$.

\textbf{19.} $-\partial_tA_z = \mu_0\dot Ir^2/4\pi R^2 = \mu_0\varepsilon_0r^2\ddot E/4$ (لأن $\dot I =
\varepsilon_0\pi R^2\ddot E$): أي تصحيح السؤال 5، حتى ثابت.

\textbf{20.} $\operatorname{div}\vect A = \partial_zA_z = 0$ لأن $A_z$ لا يتوقف إلا على $r$؛
وأي $\vect A + \operatorname{\vect{grad}}\chi$.

\textbf{21.} $\sigma = \varepsilon_0E_0 = \qty{5}{\micro C/m^2}$؛ وداخل المعدن $\vect E =
\vect 0$، ومنه تكون القفزة $\sigma/\varepsilon_0$ هي مجال الفجوة بالضبط.

\textbf{22.} يغذّي التيار $I(t)$ اللبوس من المحور؛ والجزء
الذي وراء $r$ ولم يُسلَّم بعد هو $I(1 - r^2/R^2)$، ومنه $j_s = I(1 -
r^2/R^2)/2\pi r$. وعلى الوجه الخارجي يعطي أمبير مجال السلك $\mu_0I/
2\pi r$؛ وعلى الوجه الداخلي $\mu_0Ir/2\pi R^2$؛ والفرق، $\mu_0I(1 -
r^2/R^2)/2\pi r = \mu_0j_s$، هو قفزة العبور.

\textbf{23.} اللبوسان: $R/\lambda = 3 \times 10^{-4}$، أي شبه مستقر كهربائي. والحلقة: $7 \times 10^{-4}$،
أي شبه مستقرة مغناطيسية. وعند \qty{1}{GHz}: $R/\lambda = 0.3$ --- أي جوف وهوائي.

\textbf{24.} $\lambda = \qty{3000}{km}$ من أجل حلقة بحجم المتر: \omterm{thm:b2:maxwell-equations:maxwell}{فتيار الإزاحة}
مهمَل تمامًا؛ والمجال الكهربائي الجوهري هو
مجال فاراداي، وهو ما يحفظه شبه الاستقرار المغناطيسي كاملًا.

\textbf{25.} الأول: أمبير--ماكسويل (\omterm{thm:b2:maxwell-equations:maxwell}{تيار الإزاحة})، وفاراداي من أجل
التصحيح، وكثافات الطاقة. والثاني: فاراداي المكامل، وقوة لورنتس.
والثالث: الكمونات، وبواسون (غاوس). والرابع: قفزتا العبور للمجالين $\vect E$
و $\vect B$؛ ونسبة الحجم إلى طول الموجة من أجل النظام.
\end{solution}
